% GENERATED FILE. Edit canonical lessons, then run npm run generate:books.
% canonical-source-hash: fnv1a64:c250641567790f8a
% canonical-lessons: IT-C319-carnevale, IT-C319-maschera, IT-C319-coriandoli, IT-C319-mascherarsi, IT-C319-sfilare

\chapter{Il carnevale}
\label{ch:it-il-carnevale}
\hlchaptermodality{319}

\begin{chapteropening}
\textbf{By the end of this chapter:} \emph{I can describe carnival: masks, confetti, dressing up and parading through the streets.}

Complete the last lesson of chapter 319: I can describe carnival: masks, confetti, dressing up and parading through the streets.
\end{chapteropening}

\section[il carnevale]{il carnevale — the carnival}

\label{lesson:IT-C319-carnevale}



\begin{quote}
Before the new one: say the Italian for the refreshments, then the Italian for the sugared almond.
\end{quote}

\subsection*{You'll want to know: il carnevale}
\textbf{il carnevale} — \textquotedblleft{}the carnival\textquotedblright{}.

\begin{quote}
Il carnevale finisce il martedì grasso.
\end{quote}

\begin{quote}
\emph{Carnival ends on Shrove Tuesday.}
\end{quote}

\emph{A carnevale ogni scherzo vale} — at carnival any prank goes.

\begin{cousinweb}
\emph{carnevale} is inherited from Medieval Latin \emph{carnelevāle} (EtymDB), from \emph{carō} \textquotedblleft{}meat\textquotedblright{} and \emph{levāre} \textquotedblleft{}to take away\textquotedblright{}: the farewell to meat before Lent.
\end{cousinweb}

\begin{grammarlens}[title={What you've built}]
One more word for telling, explaining and describing at B1.
\end{grammarlens}

\subsection*{Your turn}
\begin{itemize}
\raggedright
  \item \emph{Say it:} \emph{il carnevale}
  \item \emph{Say it:} \emph{il carnevale}, once more
  \item \emph{Say it:} a sentence of your own with \emph{il carnevale}
  \item \emph{Recall:} say the Italian for the refreshments, then the Italian for the sugared almond, then say \emph{il carnevale} again
\end{itemize}

\begin{tcolorbox}[breakable,colback=teal!4,colframe=teal!35!black,title={Before you move on}]
What does il carnevale mean? (\textquotedblleft{}The carnival\textquotedblright{}.) Say it once more.
\end{tcolorbox}
\section[la maschera]{la maschera — the mask}

\label{lesson:IT-C319-maschera}



\begin{quote}
Before the new one: say the Italian for the sugared almond, then the Italian for the carnival.
\end{quote}

\subsection*{You'll want to know: la maschera}
\textbf{la maschera} — \textquotedblleft{}the mask\textquotedblright{}.

\begin{quote}
Mio figlio vuole una maschera da leone.
\end{quote}

\begin{quote}
\emph{My son wants a lion mask.}
\end{quote}

Also a stock character of old comedy: \emph{le maschere della commedia}.

\begin{cousinweb}
\emph{maschera} comes from Medieval Latin \emph{masca} \textquotedblleft{}mask, disguise\textquotedblright{} (EtymDB); the deeper origin is not settled.
\end{cousinweb}

\begin{grammarlens}[title={What you've built}]
One more word for telling, explaining and describing at B1.
\end{grammarlens}

\subsection*{Your turn}
\begin{itemize}
\raggedright
  \item \emph{Say it:} \emph{la maschera}
  \item \emph{Say it:} \emph{la maschera}, once more
  \item \emph{Say it:} a sentence of your own with \emph{la maschera}
  \item \emph{Recall:} say the Italian for the sugared almond, then the Italian for the carnival, then say \emph{la maschera} again
\end{itemize}

\begin{tcolorbox}[breakable,colback=teal!4,colframe=teal!35!black,title={Before you move on}]
What does la maschera mean? (\textquotedblleft{}The mask\textquotedblright{}.) Say it once more.
\end{tcolorbox}
\section[i coriandoli]{i coriandoli — the confetti}

\label{lesson:IT-C319-coriandoli}



\begin{quote}
Before the new one: say the Italian for the carnival, then the Italian for the mask.
\end{quote}

\subsection*{You'll want to know: i coriandoli}
\textbf{i coriandoli} — \textquotedblleft{}the confetti\textquotedblright{}.

\begin{quote}
La strada è piena di coriandoli.
\end{quote}

\begin{quote}
\emph{The street is full of confetti.}
\end{quote}

Singular \emph{il coriandolo}, which is also the coriander plant.

\begin{cousinweb}
\emph{coriandoli} are named after \emph{coriandolo} \textquotedblleft{}coriander\textquotedblright{}, from Latin \emph{coriandrum} (Lewis \& Short): sugared seeds were once thrown at carnival.
\end{cousinweb}

\begin{grammarlens}[title={What you've built}]
One more word for telling, explaining and describing at B1.
\end{grammarlens}

\subsection*{Your turn}
\begin{itemize}
\raggedright
  \item \emph{Say it:} \emph{i coriandoli}
  \item \emph{Say it:} \emph{i coriandoli}, once more
  \item \emph{Say it:} a sentence of your own with \emph{i coriandoli}
  \item \emph{Recall:} say the Italian for the carnival, then the Italian for the mask, then say \emph{i coriandoli} again
\end{itemize}

\begin{tcolorbox}[breakable,colback=teal!4,colframe=teal!35!black,title={Before you move on}]
What does i coriandoli mean? (\textquotedblleft{}The confetti\textquotedblright{}.) Say it once more.
\end{tcolorbox}
\section[mascherarsi]{mascherarsi — to dress up in costume}

\label{lesson:IT-C319-mascherarsi}



\begin{quote}
Before the new one: say the Italian for the mask, then the Italian for the confetti.
\end{quote}

\subsection*{You'll want to know: mascherarsi}
\textbf{mascherarsi} — \textquotedblleft{}to dress up in costume\textquotedblright{}.

\begin{quote}
Mi sono mascherato da gatto.
\end{quote}

\begin{quote}
\emph{I dressed up as a cat.}
\end{quote}

\emph{Mascherarsi da} + a character.

\begin{cousinweb}
\emph{mascherarsi} is built on \emph{maschera} \textquotedblleft{}mask\textquotedblright{}, from Medieval Latin \emph{masca} (EtymDB).
\end{cousinweb}

\begin{grammarlens}[title={What you've built}]
One more word for telling, explaining and describing at B1.
\end{grammarlens}

\subsection*{Your turn}
\begin{itemize}
\raggedright
  \item \emph{Say it:} \emph{mascherarsi}
  \item \emph{Say it:} \emph{mascherarsi}, once more
  \item \emph{Say it:} a sentence of your own with \emph{mascherarsi}
  \item \emph{Recall:} say the Italian for the mask, then the Italian for the confetti, then say \emph{mascherarsi} again
\end{itemize}

\begin{tcolorbox}[breakable,colback=teal!4,colframe=teal!35!black,title={Before you move on}]
What does mascherarsi mean? (\textquotedblleft{}To dress up in costume\textquotedblright{}.) Say it once more.
\end{tcolorbox}
\section[sfilare]{sfilare — to parade, to march past}

\label{lesson:IT-C319-sfilare}



\begin{quote}
Before the new one: say the Italian for the confetti, then the Italian for to dress up in costume.
\end{quote}

\subsection*{You'll want to know: sfilare}
\textbf{sfilare} — \textquotedblleft{}to parade, to march past\textquotedblright{}.

\begin{quote}
I carri sfilano per le vie del centro.
\end{quote}

\begin{quote}
\emph{The floats parade through the streets of the centre.}
\end{quote}

Models \emph{sfilano} at a fashion show, too.

\begin{cousinweb}
\emph{sfilare} is built on \emph{fila} \textquotedblleft{}a line\textquotedblright{}, from Latin \emph{fīlum} \textquotedblleft{}thread\textquotedblright{} (Lewis \& Short): to go by in a line.
\end{cousinweb}

\begin{grammarlens}[title={What you've built}]
One more word for telling, explaining and describing at B1.
\end{grammarlens}

\subsection*{Your turn}
\begin{itemize}
\raggedright
  \item \emph{Say it:} \emph{sfilare}
  \item \emph{Say it:} \emph{sfilare}, once more
  \item \emph{Say it:} a sentence of your own with \emph{sfilare}
  \item \emph{Recall:} say the Italian for the confetti, then the Italian for to dress up in costume, then say \emph{sfilare} again
\end{itemize}

\begin{tcolorbox}[breakable,colback=teal!4,colframe=teal!35!black,title={Before you move on}]
What does sfilare mean? (\textquotedblleft{}To parade, to march past\textquotedblright{}.) Say it once more.
\end{tcolorbox}
